\documentclass[11pt,a4paper]{article}
% =====================================================================
%  Gemeinsame Vorlage der Arbeitsblätter – Spezielle Relativitätstheorie (Physik 11, Ph11 3.2)
%  Kopf, Fuß und Boxen wie informatik/13/rechner/arbeitsblaetter/ab-vorlage.tex,
%  Antwortplatz-Makros wie informatik/10/java/arbeitsblaetter/ab-vorlage.tex
%  (Linien 8 mm, Karo 5 mm aus ganzen Kästchen, Lücken x 1,5).
%  Einbinden in jedem Blatt direkt nach \documentclass: % =====================================================================
%  Gemeinsame Vorlage der Arbeitsblätter – Spezielle Relativitätstheorie (Physik 11, Ph11 3.2)
%  Kopf, Fuß und Boxen wie informatik/13/rechner/arbeitsblaetter/ab-vorlage.tex,
%  Antwortplatz-Makros wie informatik/10/java/arbeitsblaetter/ab-vorlage.tex
%  (Linien 8 mm, Karo 5 mm aus ganzen Kästchen, Lücken x 1,5).
%  Einbinden in jedem Blatt direkt nach \documentclass: % =====================================================================
%  Gemeinsame Vorlage der Arbeitsblätter – Spezielle Relativitätstheorie (Physik 11, Ph11 3.2)
%  Kopf, Fuß und Boxen wie informatik/13/rechner/arbeitsblaetter/ab-vorlage.tex,
%  Antwortplatz-Makros wie informatik/10/java/arbeitsblaetter/ab-vorlage.tex
%  (Linien 8 mm, Karo 5 mm aus ganzen Kästchen, Lücken x 1,5).
%  Einbinden in jedem Blatt direkt nach \documentclass: \input{ab-vorlage}
%  Übersetzen mit XeLaTeX:
%
%  Blatt:   xelatex ab-4-zeitdilatation.tex
%  Lösung:  xelatex ab-4-zeitdilatation-lsg.tex
%           (Datei enthält nur: \def\mitloesung{}\input{ab-4-zeitdilatation})
%  Lösungseinträge belegen nie mehr Platz als die Lücke im Blatt.
%
%  Antwortplatz:
%   \luecke[l] / \lsg[l]{Text}      Lücke (Standard 6,75 cm, Fachbegriff 5,25 cm, Zahl 2,25 cm)
%   \restzeile / \lsgrest{Text}     Lücke bis Zeilenende      \lueckentext  Absatz mit Lücken
%   \linien{n} / \lsglinien{n}{T}   Schreiblinien 8 mm        \schreibzeile Tabellenzelle 8 mm
%   \lsgtext{…}                     nur in der Lösung         \lsgoder{Lösung}{Blatt}
%   \kariert{n} / \lsgkariert{n}{TikZ}, \karobox{b}{h} / \lsgkarobox{b}{h}{TikZ}
%   \kreuz{1}                       Ankreuzfeld (1 = in der Lösung angekreuzt)
%   \nr{1}                          Schrittnummer ① wie auf der Website
%  Feste Farben wie auf der Website:
%   \bea{…} blau (Bahnsteig-System S), \zoe{…} orange (Zug-System S′), \licht{…} gelb
%   \weblink{Abschnitt}             grauer Verweis auf den Abschnitt der Stationsseite
% =====================================================================
\usepackage[a4paper,left=16mm,right=16mm,top=23mm,bottom=13mm,headheight=12mm,headsep=3mm,footskip=6mm]{geometry}
\usepackage{iftex}
\ifPDFTeX
  \usepackage[T1]{fontenc}
  \usepackage[utf8]{inputenc}
  \usepackage{helvet}
  \renewcommand{\familydefault}{\sfdefault}
\else
  \usepackage{fontspec}
  \setmainfont{texgyreheros}[Extension=.otf,UprightFont=*-regular,BoldFont=*-bold,ItalicFont=*-italic,BoldItalicFont=*-bolditalic]
\fi
\usepackage[ngerman,shorthands=off]{babel}
\usepackage{amsmath,amssymb}
\usepackage{xcolor,graphicx,tabularx,array,enumitem,fancyhdr,tikz,multicol,calc,etoolbox,pifont,xspace,colortbl}
\usepackage[most]{tcolorbox}
\usetikzlibrary{arrows.meta,positioning,calc}
\usepackage[hidelinks]{hyperref}
\setlength{\parindent}{0pt}
\setlength{\parskip}{3pt}
\setlist{nosep,leftmargin=*}

% ---------- Lösungsschalter ----------
\newif\ifloesung
\ifdefined\mitloesung\loesungtrue\fi
\newcommand{\lsgfarbe}{\color{red!75!black}}

% ---------- Farben wie auf der Website ----------
\definecolor{beablau}{HTML}{2563EB}\definecolor{zoeorange}{HTML}{EA580C}\definecolor{lichtgelb}{HTML}{CA8A04}
\definecolor{relgruen}{HTML}{0F766E}
\newcommand{\bea}[1]{{\color{beablau}\bfseries #1}}
\newcommand{\zoe}[1]{{\color{zoeorange}\bfseries #1}}
\newcommand{\licht}[1]{{\color{lichtgelb}\bfseries #1}}

% ---------- Kopf- und Fußzeile ----------
\newcommand{\lehrkraft}{Hau}
\newcommand{\themenbereich}{3.2 Spezielle Relativitätstheorie}
\newcommand{\abnummer}{}
\pagestyle{fancy}\fancyhf{}
\renewcommand{\headrulewidth}{0.4pt}
\fancyhead[L]{\small Klasse 11\\Datum:}
\fancyhead[C]{\small Themenbereich:\\\themenbereich}
\fancyhead[R]{\small \lehrkraft\ifloesung\\{\lsgfarbe\bfseries Lösung}\fi}
\fancyfoot[C]{\scriptsize edu-mrh.de\,·\,Physik 11\,·\,Blatt \abnummer\ifloesung\ (Lösung)\fi\,·\,Seite \thepage}
\newcommand{\abtitel}[2]{\renewcommand{\abnummer}{#1}\begin{center}\Large\bfseries #1\quad\underline{#2}\end{center}\vspace{-2mm}}
\newcommand{\leitfrage}[1]{\begin{center}\vspace{-1.5mm}\itshape #1\end{center}\vspace{-1.5mm}}

% ---------- Boxen ----------
\newenvironment{aufgabe}[2][]{\begin{tcolorbox}[enhanced,breakable,colback=white,colframe=black,boxrule=0.7pt,arc=3mm,left=2mm,right=2mm,top=1mm,bottom=1.2mm,before skip=2mm,after skip=2mm]\textbf{Aufgabe #2}\ifstrempty{#1}{}{ \textit{#1}}\par\vspace{0.3mm}}{\end{tcolorbox}}
\newtcolorbox{merke}[1][Merke]{enhanced,breakable,colback=green!5,colframe=green!40!black,boxrule=0.8pt,arc=2mm,left=2mm,right=2mm,top=1mm,bottom=1.2mm,before skip=2mm,after skip=2mm,title={#1},fonttitle=\bfseries,coltitle=white,toptitle=0.3mm,bottomtitle=0.3mm}
\newtcolorbox{hinweis}{enhanced,breakable,colback=yellow!12,colframe=orange!70!black,boxrule=0.6pt,arc=2mm,left=2mm,right=2mm,top=1mm,bottom=1mm,before skip=2mm,after skip=2mm}
\newtcolorbox{planbox}[1][Plan-Karte]{enhanced,breakable,colback=relgruen!5,colframe=relgruen,boxrule=0.9pt,arc=2mm,left=2mm,right=2mm,top=1mm,bottom=1.2mm,before skip=2mm,after skip=2mm,title={#1},fonttitle=\bfseries,coltitle=white,toptitle=0.3mm,bottomtitle=0.3mm}
\newcommand{\weblink}[1]{{\small\color{gray}\textit{Website: #1}}}

% ---------- Antwortplatz (wie informatik/10/java, Lücken x 1,5) ----------
\newlength{\zeilenabstand}\setlength{\zeilenabstand}{8mm}
\newlength{\karo}\setlength{\karo}{5mm}
\tikzset{karo/.style={draw=gray!50,line width=0.3pt},schreiblinie/.style={draw=black!65,line width=0.35pt}}
\newlength{\lsgbreite}
\newcommand{\lsg}[2][6.75cm]{\leavevmode\ifloesung\settowidth{\lsgbreite}{\,#2\,}%
  \ifdim\lsgbreite>#1\relax\rlap{\raisebox{0.8pt}{\resizebox{#1}{!}{\lsgfarbe\,#2\,}}}%
  \else\rlap{\raisebox{0.8pt}{\lsgfarbe\,#2}}\fi\fi\rule[-0.6ex]{#1}{0.4pt}\xspace}
\newcommand{\luecke}[1][6.75cm]{\lsg[#1]{}}
\newcommand{\lsgrest}[1]{\leavevmode\unskip\hspace{1.5mm}\ifloesung\rlap{\raisebox{0.8pt}{\lsgfarbe\,#1}}\fi\leaders\hbox{\rule[-0.6ex]{0.5pt}{0.4pt}}\hfill\null\par}
\newcommand{\restzeile}{\lsgrest{}}
\newcommand{\lsgtext}[1]{\ifloesung\leavevmode{\lsgfarbe#1}\fi}
\newcommand{\lsgoder}[2]{\ifloesung\leavevmode#1\else#2\fi}
\newcommand{\lueckentext}{\linespread{1.5}\selectfont\raggedright}
\newcommand{\schreibzeile}{\rule[-2.8mm]{0pt}{8mm}}
\newcommand{\nr}[1]{\ding{\numexpr171+#1\relax}}
\newcommand{\lsglinien}[2]{\par\vspace{0.5mm}\noindent\begin{tikzpicture}
  \path (0,0) (\linewidth-0.4pt,-#1\zeilenabstand-1.5mm);
  \foreach \i in {1,...,#1}{\draw[schreiblinie] (0.2pt,-\i\zeilenabstand) -- (\linewidth-0.2pt,-\i\zeilenabstand);}
  \ifloesung\node[overlay,anchor=north west,inner sep=0pt] at (1mm,-\zeilenabstand+1pt+5mm) {\parbox[t]{\linewidth-2mm}{\lsgfarbe\raggedright\setlength{\baselineskip}{\zeilenabstand}\rule{0pt}{5mm}#2}};\fi
  \end{tikzpicture}\par}
\newcommand{\linien}[1]{\lsglinien{#1}{}}
\newcommand{\karoraster}[3]{\begin{tikzpicture}\draw[karo] (0,0) grid[step=\karo] (#1\karo,#2\karo);
  \ifloesung\begin{scope}[overlay,shift={(0,#2\karo)},red!75!black,text=red!75!black,line width=0.8pt]#3\end{scope}\fi\end{tikzpicture}}
\newcommand{\lsgkariert}[2]{\par\vspace{1.5mm}\noindent\pgfmathtruncatemacro{\karoX}{floor(\linewidth/\karo)}%
  \makebox[\linewidth]{\karoraster{\karoX}{#1}{#2}}\par\vspace{1mm}}
\newcommand{\kariert}[1]{\lsgkariert{#1}{}}
\newcommand{\lsgkarobox}[3]{\pgfmathtruncatemacro{\karoX}{round((#1)/\karo)}\pgfmathtruncatemacro{\karoY}{round((#2)/\karo)}\karoraster{\karoX}{\karoY}{#3}}
\newcommand{\karobox}[2]{\lsgkarobox{#1}{#2}{}}
\newcommand{\kreuz}[1]{\leavevmode\ifloesung\ifnum#1=1 {\lsgfarbe$\boxtimes$}\else$\square$\fi\else$\square$\fi}

% ---------- Zeichnungen: Strichfigur (Füße im Ursprung des scope) ----------
%  \figur{Farbe}{Name} in einer tikzpicture, vorher mit \begin{scope}[shift={(x,y)}] platzieren
\newcommand{\figur}[2]{\draw[#1,line width=1pt,line cap=round] (0,0.62) -- (0,0.28) (0,0.28) -- (-0.12,0) (0,0.28) -- (0.12,0) (-0.15,0.48) -- (0.15,0.48);
  \fill[#1] (0,0.74) circle (0.11); \node[#1,font=\scriptsize\bfseries,above] at (0,0.82) {#2};}

%  Übersetzen mit XeLaTeX:
%
%  Blatt:   xelatex ab-4-zeitdilatation.tex
%  Lösung:  xelatex ab-4-zeitdilatation-lsg.tex
%           (Datei enthält nur: \def\mitloesung{}\documentclass[11pt,a4paper]{article}
\input{ab-vorlage}
% Blatt:   xelatex ab-4-zeitdilatation.tex
% Lösung:  xelatex ab-4-zeitdilatation-lsg.tex
% Progression: Vorhersage -> Lichtuhr, Eigen-Tick -> Messreihe + Diagramm -> Herleitung ①②③ am Dreieck
%              -> alle Uhren (Postulat 1) -> Merkkasten -> γ im Alltag -> Rechnen mit dem Plan -> Hafele-Keating
\begin{document}
\abtitel{4}{Die Lichtuhr}
\leitfrage{Gehen bewegte Uhren anders?}

\begin{aufgabe}[Vorhersage und die einfachste Uhr der Welt]{1}
\zoe{Zoes} Uhr fährt schnell an \bea{Bea} vorbei. Für Bea geht sie \ \kreuz{0} gleich schnell \ \kreuz{1} langsamer \ \kreuz{0} schneller.\\[1mm]
\lueckentext
Lichtuhr: Spiegelabstand $h = 1{,}5$ m. Ein Tick (hin und zurück) dauert für jemanden neben der Uhr\\
$\Delta t_0 = $ \lsg[8cm]{$2h/c = 3\,\text{m} : (3\cdot10^8\,\text{m/s}) = 10^{-8}\,\text{s}$} $=$ \lsg[2.25cm]{10 ns}.
\end{aufgabe}

\begin{aufgabe}[Messreihe: Zoes Lichtuhr, gemessen in Beas System]{2}
\begin{minipage}[t]{0.56\linewidth}
\vspace{0pt}
\lueckentext
\textbf{a)} Der Tick dauert für Bea länger, weil \lsgrest{das Licht einen}
\lsgrest{längeren (schrägen) Weg hat, aber auch nur mit c läuft.}
{\renewcommand{\arraystretch}{1}\small\linespread{1}\selectfont
\begin{tabularx}{\linewidth}{|l|>{\centering\arraybackslash}X|>{\centering\arraybackslash}X|>{\centering\arraybackslash}X|}\hline
$v$ & \textbf{Lichtweg in m} & \textbf{Tick in ns} & \textbf{Verhältnis}\\\hline
0\schreibzeile & \lsgtext{3,00} & \lsgtext{10,00} & \lsgtext{1,00}\\\hline
0,3 c\schreibzeile & \lsgtext{3,14} & \lsgtext{10,48} & \lsgtext{1,05}\\\hline
0,6 c\schreibzeile & \lsgtext{3,75} & \lsgtext{12,50} & \lsgtext{1,25}\\\hline
0,8 c\schreibzeile & \lsgtext{5,00} & \lsgtext{16,67} & \lsgtext{1,67}\\\hline
0,9 c\schreibzeile & \lsgtext{6,88} & \lsgtext{22,94} & \lsgtext{2,29}\\\hline
\end{tabularx}}

\vspace{1mm}
\textbf{d)} Beas Uhr zeigt 10 Ticks, Zoes Uhr bei 0,6 c: \lsg[1.5cm]{8}
\end{minipage}\hfill
\begin{minipage}[t]{0.41\linewidth}
\vspace{0pt}\centering
{\small \textbf{c)} Verhältnis (hoch) gegen $v$ (rechts)}\\[0.5mm]
\begin{tikzpicture}[x=5mm,y=5mm]
  \draw[karo] (0,0) grid (12,12);
  \draw[-{Stealth},line width=0.7pt] (1,1) -- (11.7,1) node[below left,font=\scriptsize] {$v$};
  \draw[-{Stealth},line width=0.7pt] (1,1) -- (1,11.7);
  \foreach \x/\l in {1/0,3/{0,2},5/{0,4},7/{0,6},9/{0,8},11/{c}} \node[below,font=\tiny] at (\x,1) {\l};
  \foreach \y/\l in {1/{1,0},5/{1,5},9/{2,0}} \node[left,font=\tiny] at (1,\y) {\l};
  \ifloesung\draw[red!75!black,line width=0.9pt] plot[smooth] coordinates {(1,1) (4,1.39) (7,3) (9,6.33) (10,11.35)};
  \foreach \p in {(1,1),(4,1.39),(7,3),(9,6.33),(10,11.35)} \fill[red!75!black] \p circle (1.6pt);\fi
\end{tikzpicture}\\[0.5mm]
\raggedright\small Verlauf: \lsgtext{erst kaum, dann immer steiler}
\end{minipage}
\end{aufgabe}

\begin{aufgabe}[Herleitung: Das Dreieck kommt zurück]{3}
\begin{minipage}[t]{0.4\linewidth}
\vspace{0pt}\centering
\begin{tikzpicture}[x=1cm,y=1cm,font=\small]
  \draw[beablau,line width=1.4pt] (0,0) -- (3,0); \draw[beablau,line width=0.6pt,dashed] (3,0) -- (6,0);
  \draw[zoeorange,line width=1.4pt] (3,0) -- (3,2.4);
  \draw[lichtgelb,line width=1.4pt] (0,0) -- (3,2.4); \draw[lichtgelb,line width=0.6pt,dashed] (3,2.4) -- (6,0);
  \draw (2.75,0) |- (3,0.25);
  \node[lichtgelb,above,rotate=38.7] at (1.4,1.2) {\lsgoder{$c\cdot\Delta t/2$}{\rule{1.6cm}{0.4pt}}};
  \node[beablau,below] at (1.5,0) {\lsgoder{$v\cdot\Delta t/2$}{\rule[-3mm]{1.6cm}{0.4pt}}};
  \node[zoeorange,right] at (3,1.1) {\lsgoder{$c\cdot\Delta t_0/2$}{\rule{1.6cm}{0.4pt}}};
  \node[font=\scriptsize,below] at (6,0) {Tick};
\end{tikzpicture}\\
{\scriptsize Ein Tick von Zoes Uhr in Beas System}
\end{minipage}\hfill
\begin{minipage}[t]{0.57\linewidth}
\vspace{0pt}\lueckentext
Derselbe Tick dauert für \bea{Bea} $\Delta t$, für \zoe{Zoe} $\Delta t_0$.\\
\nr{1} \textbf{Seiten} beschriften (halbe Tick-Dauer, im Bild).\\
\nr{2} \textbf{Pythagoras:}\\ \lsgrest{$\left(c\,\frac{\Delta t}{2}\right)^2 = \left(v\,\frac{\Delta t}{2}\right)^2 + \left(c\,\frac{\Delta t_0}{2}\right)^2$}
\nr{3} \textbf{Auflösen} nach $\Delta t$:\\ \lsgrest{$c^2\Delta t^2 - v^2\Delta t^2 = c^2\Delta t_0^2$}
\lsgrest{$\Delta t^2\,(1 - v^2/c^2) = \Delta t_0^2$}
$\Delta t = $ \lsg[5cm]{$\Delta t_0 / \sqrt{1 - v^2/c^2}$}
\end{minipage}\\[1mm]
\lueckentext
Prüfen mit $v = 0{,}6\,c$: \lsgrest{$\sqrt{1 - 0{,}36} = 0{,}8$; \quad $\Delta t = 10$ ns : 0,8 = 12,5 ns \checkmark}
\end{aufgabe}

\newpage
\begin{aufgabe}[Gilt das nur für Lichtuhren? \normalfont(Argumentiere mit Postulat 1)]{4}
\lsglinien{3}{Nein. Ginge Zoes Armbanduhr anders als ihre Lichtuhr, könnte sie im geschlossenen Zug feststellen, dass sie fährt. Das verbietet Postulat 1. Also laufen alle Uhren (auch Puls, Zerfälle) im Gleichschritt.}
\end{aufgabe}

\begin{merke}[Zeitdilatation]
\lueckentext
Die \textbf{Eigenzeit} $\Delta t_0$ misst, für wen beide Ereignisse \lsg[5.25cm]{am selben Ort} stattfinden.\\
In jedem anderen Inertialsystem liegt zwischen den Ereignissen eine \lsg[3.2cm]{längere} Zeit:\\
$\Delta t = $ \lsg[3cm]{$\gamma\cdot\Delta t_0$} \quad mit dem Lorentzfaktor \quad $\gamma = $ \lsg[4.5cm]{$1/\sqrt{1 - v^2/c^2}$}\\
$\gamma$ ist immer \lsg[1.5cm]{$\geq 1$}. Bewegte Uhren gehen \lsg[3.2cm]{langsamer}. Die Eigenzeit ist die \lsg[3.2cm]{kürzeste} Zeit.
\end{merke}

\begin{aufgabe}[Warum hat das vorher niemand bemerkt? \normalfont($\gamma$-Explorer)]{5}
{\renewcommand{\arraystretch}{1}\small
\begin{tabularx}{\linewidth}{|l|>{\centering\arraybackslash}X|>{\centering\arraybackslash}X|>{\centering\arraybackslash}X|>{\centering\arraybackslash}X|}\hline
 & \textbf{ICE} & \textbf{Flugzeug} & \textbf{ISS} & \textbf{Myon (0,995 c)}\\\hline
$\gamma$ bzw. $\gamma - 1$\schreibzeile & \lsgtext{$1 + 3{,}8\cdot10^{-14}$} & \lsgtext{$1 + 3{,}5\cdot10^{-13}$} & \lsgtext{$1 + 3{,}3\cdot10^{-10}$} & \lsgtext{10,0}\\\hline
Uhr geht nach pro Tag\schreibzeile & \lsgtext{3,3 ns} & \lsgtext{30 ns} & \lsgtext{28 µs} & \lsgtext{(über 21 h)}\\\hline
\end{tabularx}}

\vspace{1mm}\lueckentext
Mein Tipp (ICE): \luecke[2.6cm] \quad $\gamma$ weicht erst ab etwa \lsg[2cm]{0,14 c} um mehr als 1\,\% von 1 ab.
\end{aufgabe}

\begin{aufgabe}[Mit dem Plan rechnen: Raumschiff mit $v = 0{,}8\,c$, Borduhr zeigt 3,0 Jahre]{6}
\lueckentext
\nr{1} \bea{S:} \lsg[4cm]{Erde und Station} \zoe{S$'$:} \lsg[3cm]{Raumschiff} Ereignisse: \lsg[4.6cm]{Abflug und Ankunft}\\
\nr{2} Beide Ereignisse am selben Ort für \lsg[3.5cm]{das Raumschiff}: Eigenzeit misst die \lsg[3cm]{Borduhr}.\\
\nr{3} $\gamma = $ \lsg[5.5cm]{$1/\sqrt{1 - 0{,}64} = 1/0{,}6 = 1{,}67$} \quad $\Delta t = $ \lsg[4.5cm]{$1{,}67\cdot 3{,}0$ a $= 5{,}0$ a}\\
\phantom{\nr{3}} Entfernung der Station für die Erde: \lsgrest{$0{,}8\,c\cdot 5{,}0$ a $= 4{,}0$ Lichtjahre}
\end{aufgabe}

\begin{aufgabe}[Im Experiment geprüft: Atomuhren im Flugzeug (1971)]{7}
{\small Vorhersage / Messung: ostwärts $-40 \pm 23$ ns / $-59 \pm 10$ ns; \quad westwärts $+275 \pm 21$ ns / $+273 \pm 7$ ns.}\\[0.5mm]
\lueckentext
\textbf{a)} Passt das zusammen? \lsgrest{Ja, die Bereiche überlappen jeweils.}
\textbf{b)} Warum geht die Uhr im Ostflug nach, die im Westflug vor?
\lsglinien{2}{Von außen gesehen ist die Uhr im Ostflug am schnellsten (Erddrehung plus Flug), die im Westflug am langsamsten. Die schnellste Uhr geht am langsamsten.}
\end{aufgabe}

\begin{planbox}[Plan-Karte, Stand nach Station 4 \normalfont(ins Logbuch übertragen, Annahmen-Tafel Spalte 4)]
\lueckentext
\textbf{\nr{2} Eigen-System:} beide Ereignisse am \lsg[3cm]{selben Ort} $\to$ Eigenzeit. \quad \textbf{\nr{3} Umrechnen:} Zeit \lsg[1.2cm]{mal} $\gamma$.
\end{planbox}
\end{document}
)
%  Lösungseinträge belegen nie mehr Platz als die Lücke im Blatt.
%
%  Antwortplatz:
%   \luecke[l] / \lsg[l]{Text}      Lücke (Standard 6,75 cm, Fachbegriff 5,25 cm, Zahl 2,25 cm)
%   \restzeile / \lsgrest{Text}     Lücke bis Zeilenende      \lueckentext  Absatz mit Lücken
%   \linien{n} / \lsglinien{n}{T}   Schreiblinien 8 mm        \schreibzeile Tabellenzelle 8 mm
%   \lsgtext{…}                     nur in der Lösung         \lsgoder{Lösung}{Blatt}
%   \kariert{n} / \lsgkariert{n}{TikZ}, \karobox{b}{h} / \lsgkarobox{b}{h}{TikZ}
%   \kreuz{1}                       Ankreuzfeld (1 = in der Lösung angekreuzt)
%   \nr{1}                          Schrittnummer ① wie auf der Website
%  Feste Farben wie auf der Website:
%   \bea{…} blau (Bahnsteig-System S), \zoe{…} orange (Zug-System S′), \licht{…} gelb
%   \weblink{Abschnitt}             grauer Verweis auf den Abschnitt der Stationsseite
% =====================================================================
\usepackage[a4paper,left=16mm,right=16mm,top=23mm,bottom=13mm,headheight=12mm,headsep=3mm,footskip=6mm]{geometry}
\usepackage{iftex}
\ifPDFTeX
  \usepackage[T1]{fontenc}
  \usepackage[utf8]{inputenc}
  \usepackage{helvet}
  \renewcommand{\familydefault}{\sfdefault}
\else
  \usepackage{fontspec}
  \setmainfont{texgyreheros}[Extension=.otf,UprightFont=*-regular,BoldFont=*-bold,ItalicFont=*-italic,BoldItalicFont=*-bolditalic]
\fi
\usepackage[ngerman,shorthands=off]{babel}
\usepackage{amsmath,amssymb}
\usepackage{xcolor,graphicx,tabularx,array,enumitem,fancyhdr,tikz,multicol,calc,etoolbox,pifont,xspace,colortbl}
\usepackage[most]{tcolorbox}
\usetikzlibrary{arrows.meta,positioning,calc}
\usepackage[hidelinks]{hyperref}
\setlength{\parindent}{0pt}
\setlength{\parskip}{3pt}
\setlist{nosep,leftmargin=*}

% ---------- Lösungsschalter ----------
\newif\ifloesung
\ifdefined\mitloesung\loesungtrue\fi
\newcommand{\lsgfarbe}{\color{red!75!black}}

% ---------- Farben wie auf der Website ----------
\definecolor{beablau}{HTML}{2563EB}\definecolor{zoeorange}{HTML}{EA580C}\definecolor{lichtgelb}{HTML}{CA8A04}
\definecolor{relgruen}{HTML}{0F766E}
\newcommand{\bea}[1]{{\color{beablau}\bfseries #1}}
\newcommand{\zoe}[1]{{\color{zoeorange}\bfseries #1}}
\newcommand{\licht}[1]{{\color{lichtgelb}\bfseries #1}}

% ---------- Kopf- und Fußzeile ----------
\newcommand{\lehrkraft}{Hau}
\newcommand{\themenbereich}{3.2 Spezielle Relativitätstheorie}
\newcommand{\abnummer}{}
\pagestyle{fancy}\fancyhf{}
\renewcommand{\headrulewidth}{0.4pt}
\fancyhead[L]{\small Klasse 11\\Datum:}
\fancyhead[C]{\small Themenbereich:\\\themenbereich}
\fancyhead[R]{\small \lehrkraft\ifloesung\\{\lsgfarbe\bfseries Lösung}\fi}
\fancyfoot[C]{\scriptsize edu-mrh.de\,·\,Physik 11\,·\,Blatt \abnummer\ifloesung\ (Lösung)\fi\,·\,Seite \thepage}
\newcommand{\abtitel}[2]{\renewcommand{\abnummer}{#1}\begin{center}\Large\bfseries #1\quad\underline{#2}\end{center}\vspace{-2mm}}
\newcommand{\leitfrage}[1]{\begin{center}\vspace{-1.5mm}\itshape #1\end{center}\vspace{-1.5mm}}

% ---------- Boxen ----------
\newenvironment{aufgabe}[2][]{\begin{tcolorbox}[enhanced,breakable,colback=white,colframe=black,boxrule=0.7pt,arc=3mm,left=2mm,right=2mm,top=1mm,bottom=1.2mm,before skip=2mm,after skip=2mm]\textbf{Aufgabe #2}\ifstrempty{#1}{}{ \textit{#1}}\par\vspace{0.3mm}}{\end{tcolorbox}}
\newtcolorbox{merke}[1][Merke]{enhanced,breakable,colback=green!5,colframe=green!40!black,boxrule=0.8pt,arc=2mm,left=2mm,right=2mm,top=1mm,bottom=1.2mm,before skip=2mm,after skip=2mm,title={#1},fonttitle=\bfseries,coltitle=white,toptitle=0.3mm,bottomtitle=0.3mm}
\newtcolorbox{hinweis}{enhanced,breakable,colback=yellow!12,colframe=orange!70!black,boxrule=0.6pt,arc=2mm,left=2mm,right=2mm,top=1mm,bottom=1mm,before skip=2mm,after skip=2mm}
\newtcolorbox{planbox}[1][Plan-Karte]{enhanced,breakable,colback=relgruen!5,colframe=relgruen,boxrule=0.9pt,arc=2mm,left=2mm,right=2mm,top=1mm,bottom=1.2mm,before skip=2mm,after skip=2mm,title={#1},fonttitle=\bfseries,coltitle=white,toptitle=0.3mm,bottomtitle=0.3mm}
\newcommand{\weblink}[1]{{\small\color{gray}\textit{Website: #1}}}

% ---------- Antwortplatz (wie informatik/10/java, Lücken x 1,5) ----------
\newlength{\zeilenabstand}\setlength{\zeilenabstand}{8mm}
\newlength{\karo}\setlength{\karo}{5mm}
\tikzset{karo/.style={draw=gray!50,line width=0.3pt},schreiblinie/.style={draw=black!65,line width=0.35pt}}
\newlength{\lsgbreite}
\newcommand{\lsg}[2][6.75cm]{\leavevmode\ifloesung\settowidth{\lsgbreite}{\,#2\,}%
  \ifdim\lsgbreite>#1\relax\rlap{\raisebox{0.8pt}{\resizebox{#1}{!}{\lsgfarbe\,#2\,}}}%
  \else\rlap{\raisebox{0.8pt}{\lsgfarbe\,#2}}\fi\fi\rule[-0.6ex]{#1}{0.4pt}\xspace}
\newcommand{\luecke}[1][6.75cm]{\lsg[#1]{}}
\newcommand{\lsgrest}[1]{\leavevmode\unskip\hspace{1.5mm}\ifloesung\rlap{\raisebox{0.8pt}{\lsgfarbe\,#1}}\fi\leaders\hbox{\rule[-0.6ex]{0.5pt}{0.4pt}}\hfill\null\par}
\newcommand{\restzeile}{\lsgrest{}}
\newcommand{\lsgtext}[1]{\ifloesung\leavevmode{\lsgfarbe#1}\fi}
\newcommand{\lsgoder}[2]{\ifloesung\leavevmode#1\else#2\fi}
\newcommand{\lueckentext}{\linespread{1.5}\selectfont\raggedright}
\newcommand{\schreibzeile}{\rule[-2.8mm]{0pt}{8mm}}
\newcommand{\nr}[1]{\ding{\numexpr171+#1\relax}}
\newcommand{\lsglinien}[2]{\par\vspace{0.5mm}\noindent\begin{tikzpicture}
  \path (0,0) (\linewidth-0.4pt,-#1\zeilenabstand-1.5mm);
  \foreach \i in {1,...,#1}{\draw[schreiblinie] (0.2pt,-\i\zeilenabstand) -- (\linewidth-0.2pt,-\i\zeilenabstand);}
  \ifloesung\node[overlay,anchor=north west,inner sep=0pt] at (1mm,-\zeilenabstand+1pt+5mm) {\parbox[t]{\linewidth-2mm}{\lsgfarbe\raggedright\setlength{\baselineskip}{\zeilenabstand}\rule{0pt}{5mm}#2}};\fi
  \end{tikzpicture}\par}
\newcommand{\linien}[1]{\lsglinien{#1}{}}
\newcommand{\karoraster}[3]{\begin{tikzpicture}\draw[karo] (0,0) grid[step=\karo] (#1\karo,#2\karo);
  \ifloesung\begin{scope}[overlay,shift={(0,#2\karo)},red!75!black,text=red!75!black,line width=0.8pt]#3\end{scope}\fi\end{tikzpicture}}
\newcommand{\lsgkariert}[2]{\par\vspace{1.5mm}\noindent\pgfmathtruncatemacro{\karoX}{floor(\linewidth/\karo)}%
  \makebox[\linewidth]{\karoraster{\karoX}{#1}{#2}}\par\vspace{1mm}}
\newcommand{\kariert}[1]{\lsgkariert{#1}{}}
\newcommand{\lsgkarobox}[3]{\pgfmathtruncatemacro{\karoX}{round((#1)/\karo)}\pgfmathtruncatemacro{\karoY}{round((#2)/\karo)}\karoraster{\karoX}{\karoY}{#3}}
\newcommand{\karobox}[2]{\lsgkarobox{#1}{#2}{}}
\newcommand{\kreuz}[1]{\leavevmode\ifloesung\ifnum#1=1 {\lsgfarbe$\boxtimes$}\else$\square$\fi\else$\square$\fi}

% ---------- Zeichnungen: Strichfigur (Füße im Ursprung des scope) ----------
%  \figur{Farbe}{Name} in einer tikzpicture, vorher mit \begin{scope}[shift={(x,y)}] platzieren
\newcommand{\figur}[2]{\draw[#1,line width=1pt,line cap=round] (0,0.62) -- (0,0.28) (0,0.28) -- (-0.12,0) (0,0.28) -- (0.12,0) (-0.15,0.48) -- (0.15,0.48);
  \fill[#1] (0,0.74) circle (0.11); \node[#1,font=\scriptsize\bfseries,above] at (0,0.82) {#2};}

%  Übersetzen mit XeLaTeX:
%
%  Blatt:   xelatex ab-4-zeitdilatation.tex
%  Lösung:  xelatex ab-4-zeitdilatation-lsg.tex
%           (Datei enthält nur: \def\mitloesung{}\documentclass[11pt,a4paper]{article}
% =====================================================================
%  Gemeinsame Vorlage der Arbeitsblätter – Spezielle Relativitätstheorie (Physik 11, Ph11 3.2)
%  Kopf, Fuß und Boxen wie informatik/13/rechner/arbeitsblaetter/ab-vorlage.tex,
%  Antwortplatz-Makros wie informatik/10/java/arbeitsblaetter/ab-vorlage.tex
%  (Linien 8 mm, Karo 5 mm aus ganzen Kästchen, Lücken x 1,5).
%  Einbinden in jedem Blatt direkt nach \documentclass: \input{ab-vorlage}
%  Übersetzen mit XeLaTeX:
%
%  Blatt:   xelatex ab-4-zeitdilatation.tex
%  Lösung:  xelatex ab-4-zeitdilatation-lsg.tex
%           (Datei enthält nur: \def\mitloesung{}\input{ab-4-zeitdilatation})
%  Lösungseinträge belegen nie mehr Platz als die Lücke im Blatt.
%
%  Antwortplatz:
%   \luecke[l] / \lsg[l]{Text}      Lücke (Standard 6,75 cm, Fachbegriff 5,25 cm, Zahl 2,25 cm)
%   \restzeile / \lsgrest{Text}     Lücke bis Zeilenende      \lueckentext  Absatz mit Lücken
%   \linien{n} / \lsglinien{n}{T}   Schreiblinien 8 mm        \schreibzeile Tabellenzelle 8 mm
%   \lsgtext{…}                     nur in der Lösung         \lsgoder{Lösung}{Blatt}
%   \kariert{n} / \lsgkariert{n}{TikZ}, \karobox{b}{h} / \lsgkarobox{b}{h}{TikZ}
%   \kreuz{1}                       Ankreuzfeld (1 = in der Lösung angekreuzt)
%   \nr{1}                          Schrittnummer ① wie auf der Website
%  Feste Farben wie auf der Website:
%   \bea{…} blau (Bahnsteig-System S), \zoe{…} orange (Zug-System S′), \licht{…} gelb
%   \weblink{Abschnitt}             grauer Verweis auf den Abschnitt der Stationsseite
% =====================================================================
\usepackage[a4paper,left=16mm,right=16mm,top=23mm,bottom=13mm,headheight=12mm,headsep=3mm,footskip=6mm]{geometry}
\usepackage{iftex}
\ifPDFTeX
  \usepackage[T1]{fontenc}
  \usepackage[utf8]{inputenc}
  \usepackage{helvet}
  \renewcommand{\familydefault}{\sfdefault}
\else
  \usepackage{fontspec}
  \setmainfont{texgyreheros}[Extension=.otf,UprightFont=*-regular,BoldFont=*-bold,ItalicFont=*-italic,BoldItalicFont=*-bolditalic]
\fi
\usepackage[ngerman,shorthands=off]{babel}
\usepackage{amsmath,amssymb}
\usepackage{xcolor,graphicx,tabularx,array,enumitem,fancyhdr,tikz,multicol,calc,etoolbox,pifont,xspace,colortbl}
\usepackage[most]{tcolorbox}
\usetikzlibrary{arrows.meta,positioning,calc}
\usepackage[hidelinks]{hyperref}
\setlength{\parindent}{0pt}
\setlength{\parskip}{3pt}
\setlist{nosep,leftmargin=*}

% ---------- Lösungsschalter ----------
\newif\ifloesung
\ifdefined\mitloesung\loesungtrue\fi
\newcommand{\lsgfarbe}{\color{red!75!black}}

% ---------- Farben wie auf der Website ----------
\definecolor{beablau}{HTML}{2563EB}\definecolor{zoeorange}{HTML}{EA580C}\definecolor{lichtgelb}{HTML}{CA8A04}
\definecolor{relgruen}{HTML}{0F766E}
\newcommand{\bea}[1]{{\color{beablau}\bfseries #1}}
\newcommand{\zoe}[1]{{\color{zoeorange}\bfseries #1}}
\newcommand{\licht}[1]{{\color{lichtgelb}\bfseries #1}}

% ---------- Kopf- und Fußzeile ----------
\newcommand{\lehrkraft}{Hau}
\newcommand{\themenbereich}{3.2 Spezielle Relativitätstheorie}
\newcommand{\abnummer}{}
\pagestyle{fancy}\fancyhf{}
\renewcommand{\headrulewidth}{0.4pt}
\fancyhead[L]{\small Klasse 11\\Datum:}
\fancyhead[C]{\small Themenbereich:\\\themenbereich}
\fancyhead[R]{\small \lehrkraft\ifloesung\\{\lsgfarbe\bfseries Lösung}\fi}
\fancyfoot[C]{\scriptsize edu-mrh.de\,·\,Physik 11\,·\,Blatt \abnummer\ifloesung\ (Lösung)\fi\,·\,Seite \thepage}
\newcommand{\abtitel}[2]{\renewcommand{\abnummer}{#1}\begin{center}\Large\bfseries #1\quad\underline{#2}\end{center}\vspace{-2mm}}
\newcommand{\leitfrage}[1]{\begin{center}\vspace{-1.5mm}\itshape #1\end{center}\vspace{-1.5mm}}

% ---------- Boxen ----------
\newenvironment{aufgabe}[2][]{\begin{tcolorbox}[enhanced,breakable,colback=white,colframe=black,boxrule=0.7pt,arc=3mm,left=2mm,right=2mm,top=1mm,bottom=1.2mm,before skip=2mm,after skip=2mm]\textbf{Aufgabe #2}\ifstrempty{#1}{}{ \textit{#1}}\par\vspace{0.3mm}}{\end{tcolorbox}}
\newtcolorbox{merke}[1][Merke]{enhanced,breakable,colback=green!5,colframe=green!40!black,boxrule=0.8pt,arc=2mm,left=2mm,right=2mm,top=1mm,bottom=1.2mm,before skip=2mm,after skip=2mm,title={#1},fonttitle=\bfseries,coltitle=white,toptitle=0.3mm,bottomtitle=0.3mm}
\newtcolorbox{hinweis}{enhanced,breakable,colback=yellow!12,colframe=orange!70!black,boxrule=0.6pt,arc=2mm,left=2mm,right=2mm,top=1mm,bottom=1mm,before skip=2mm,after skip=2mm}
\newtcolorbox{planbox}[1][Plan-Karte]{enhanced,breakable,colback=relgruen!5,colframe=relgruen,boxrule=0.9pt,arc=2mm,left=2mm,right=2mm,top=1mm,bottom=1.2mm,before skip=2mm,after skip=2mm,title={#1},fonttitle=\bfseries,coltitle=white,toptitle=0.3mm,bottomtitle=0.3mm}
\newcommand{\weblink}[1]{{\small\color{gray}\textit{Website: #1}}}

% ---------- Antwortplatz (wie informatik/10/java, Lücken x 1,5) ----------
\newlength{\zeilenabstand}\setlength{\zeilenabstand}{8mm}
\newlength{\karo}\setlength{\karo}{5mm}
\tikzset{karo/.style={draw=gray!50,line width=0.3pt},schreiblinie/.style={draw=black!65,line width=0.35pt}}
\newlength{\lsgbreite}
\newcommand{\lsg}[2][6.75cm]{\leavevmode\ifloesung\settowidth{\lsgbreite}{\,#2\,}%
  \ifdim\lsgbreite>#1\relax\rlap{\raisebox{0.8pt}{\resizebox{#1}{!}{\lsgfarbe\,#2\,}}}%
  \else\rlap{\raisebox{0.8pt}{\lsgfarbe\,#2}}\fi\fi\rule[-0.6ex]{#1}{0.4pt}\xspace}
\newcommand{\luecke}[1][6.75cm]{\lsg[#1]{}}
\newcommand{\lsgrest}[1]{\leavevmode\unskip\hspace{1.5mm}\ifloesung\rlap{\raisebox{0.8pt}{\lsgfarbe\,#1}}\fi\leaders\hbox{\rule[-0.6ex]{0.5pt}{0.4pt}}\hfill\null\par}
\newcommand{\restzeile}{\lsgrest{}}
\newcommand{\lsgtext}[1]{\ifloesung\leavevmode{\lsgfarbe#1}\fi}
\newcommand{\lsgoder}[2]{\ifloesung\leavevmode#1\else#2\fi}
\newcommand{\lueckentext}{\linespread{1.5}\selectfont\raggedright}
\newcommand{\schreibzeile}{\rule[-2.8mm]{0pt}{8mm}}
\newcommand{\nr}[1]{\ding{\numexpr171+#1\relax}}
\newcommand{\lsglinien}[2]{\par\vspace{0.5mm}\noindent\begin{tikzpicture}
  \path (0,0) (\linewidth-0.4pt,-#1\zeilenabstand-1.5mm);
  \foreach \i in {1,...,#1}{\draw[schreiblinie] (0.2pt,-\i\zeilenabstand) -- (\linewidth-0.2pt,-\i\zeilenabstand);}
  \ifloesung\node[overlay,anchor=north west,inner sep=0pt] at (1mm,-\zeilenabstand+1pt+5mm) {\parbox[t]{\linewidth-2mm}{\lsgfarbe\raggedright\setlength{\baselineskip}{\zeilenabstand}\rule{0pt}{5mm}#2}};\fi
  \end{tikzpicture}\par}
\newcommand{\linien}[1]{\lsglinien{#1}{}}
\newcommand{\karoraster}[3]{\begin{tikzpicture}\draw[karo] (0,0) grid[step=\karo] (#1\karo,#2\karo);
  \ifloesung\begin{scope}[overlay,shift={(0,#2\karo)},red!75!black,text=red!75!black,line width=0.8pt]#3\end{scope}\fi\end{tikzpicture}}
\newcommand{\lsgkariert}[2]{\par\vspace{1.5mm}\noindent\pgfmathtruncatemacro{\karoX}{floor(\linewidth/\karo)}%
  \makebox[\linewidth]{\karoraster{\karoX}{#1}{#2}}\par\vspace{1mm}}
\newcommand{\kariert}[1]{\lsgkariert{#1}{}}
\newcommand{\lsgkarobox}[3]{\pgfmathtruncatemacro{\karoX}{round((#1)/\karo)}\pgfmathtruncatemacro{\karoY}{round((#2)/\karo)}\karoraster{\karoX}{\karoY}{#3}}
\newcommand{\karobox}[2]{\lsgkarobox{#1}{#2}{}}
\newcommand{\kreuz}[1]{\leavevmode\ifloesung\ifnum#1=1 {\lsgfarbe$\boxtimes$}\else$\square$\fi\else$\square$\fi}

% ---------- Zeichnungen: Strichfigur (Füße im Ursprung des scope) ----------
%  \figur{Farbe}{Name} in einer tikzpicture, vorher mit \begin{scope}[shift={(x,y)}] platzieren
\newcommand{\figur}[2]{\draw[#1,line width=1pt,line cap=round] (0,0.62) -- (0,0.28) (0,0.28) -- (-0.12,0) (0,0.28) -- (0.12,0) (-0.15,0.48) -- (0.15,0.48);
  \fill[#1] (0,0.74) circle (0.11); \node[#1,font=\scriptsize\bfseries,above] at (0,0.82) {#2};}

% Blatt:   xelatex ab-4-zeitdilatation.tex
% Lösung:  xelatex ab-4-zeitdilatation-lsg.tex
% Progression: Vorhersage -> Lichtuhr, Eigen-Tick -> Messreihe + Diagramm -> Herleitung ①②③ am Dreieck
%              -> alle Uhren (Postulat 1) -> Merkkasten -> γ im Alltag -> Rechnen mit dem Plan -> Hafele-Keating
\begin{document}
\abtitel{4}{Die Lichtuhr}
\leitfrage{Gehen bewegte Uhren anders?}

\begin{aufgabe}[Vorhersage und die einfachste Uhr der Welt]{1}
\zoe{Zoes} Uhr fährt schnell an \bea{Bea} vorbei. Für Bea geht sie \ \kreuz{0} gleich schnell \ \kreuz{1} langsamer \ \kreuz{0} schneller.\\[1mm]
\lueckentext
Lichtuhr: Spiegelabstand $h = 1{,}5$ m. Ein Tick (hin und zurück) dauert für jemanden neben der Uhr\\
$\Delta t_0 = $ \lsg[8cm]{$2h/c = 3\,\text{m} : (3\cdot10^8\,\text{m/s}) = 10^{-8}\,\text{s}$} $=$ \lsg[2.25cm]{10 ns}.
\end{aufgabe}

\begin{aufgabe}[Messreihe: Zoes Lichtuhr, gemessen in Beas System]{2}
\begin{minipage}[t]{0.56\linewidth}
\vspace{0pt}
\lueckentext
\textbf{a)} Der Tick dauert für Bea länger, weil \lsgrest{das Licht einen}
\lsgrest{längeren (schrägen) Weg hat, aber auch nur mit c läuft.}
{\renewcommand{\arraystretch}{1}\small\linespread{1}\selectfont
\begin{tabularx}{\linewidth}{|l|>{\centering\arraybackslash}X|>{\centering\arraybackslash}X|>{\centering\arraybackslash}X|}\hline
$v$ & \textbf{Lichtweg in m} & \textbf{Tick in ns} & \textbf{Verhältnis}\\\hline
0\schreibzeile & \lsgtext{3,00} & \lsgtext{10,00} & \lsgtext{1,00}\\\hline
0,3 c\schreibzeile & \lsgtext{3,14} & \lsgtext{10,48} & \lsgtext{1,05}\\\hline
0,6 c\schreibzeile & \lsgtext{3,75} & \lsgtext{12,50} & \lsgtext{1,25}\\\hline
0,8 c\schreibzeile & \lsgtext{5,00} & \lsgtext{16,67} & \lsgtext{1,67}\\\hline
0,9 c\schreibzeile & \lsgtext{6,88} & \lsgtext{22,94} & \lsgtext{2,29}\\\hline
\end{tabularx}}

\vspace{1mm}
\textbf{d)} Beas Uhr zeigt 10 Ticks, Zoes Uhr bei 0,6 c: \lsg[1.5cm]{8}
\end{minipage}\hfill
\begin{minipage}[t]{0.41\linewidth}
\vspace{0pt}\centering
{\small \textbf{c)} Verhältnis (hoch) gegen $v$ (rechts)}\\[0.5mm]
\begin{tikzpicture}[x=5mm,y=5mm]
  \draw[karo] (0,0) grid (12,12);
  \draw[-{Stealth},line width=0.7pt] (1,1) -- (11.7,1) node[below left,font=\scriptsize] {$v$};
  \draw[-{Stealth},line width=0.7pt] (1,1) -- (1,11.7);
  \foreach \x/\l in {1/0,3/{0,2},5/{0,4},7/{0,6},9/{0,8},11/{c}} \node[below,font=\tiny] at (\x,1) {\l};
  \foreach \y/\l in {1/{1,0},5/{1,5},9/{2,0}} \node[left,font=\tiny] at (1,\y) {\l};
  \ifloesung\draw[red!75!black,line width=0.9pt] plot[smooth] coordinates {(1,1) (4,1.39) (7,3) (9,6.33) (10,11.35)};
  \foreach \p in {(1,1),(4,1.39),(7,3),(9,6.33),(10,11.35)} \fill[red!75!black] \p circle (1.6pt);\fi
\end{tikzpicture}\\[0.5mm]
\raggedright\small Verlauf: \lsgtext{erst kaum, dann immer steiler}
\end{minipage}
\end{aufgabe}

\begin{aufgabe}[Herleitung: Das Dreieck kommt zurück]{3}
\begin{minipage}[t]{0.4\linewidth}
\vspace{0pt}\centering
\begin{tikzpicture}[x=1cm,y=1cm,font=\small]
  \draw[beablau,line width=1.4pt] (0,0) -- (3,0); \draw[beablau,line width=0.6pt,dashed] (3,0) -- (6,0);
  \draw[zoeorange,line width=1.4pt] (3,0) -- (3,2.4);
  \draw[lichtgelb,line width=1.4pt] (0,0) -- (3,2.4); \draw[lichtgelb,line width=0.6pt,dashed] (3,2.4) -- (6,0);
  \draw (2.75,0) |- (3,0.25);
  \node[lichtgelb,above,rotate=38.7] at (1.4,1.2) {\lsgoder{$c\cdot\Delta t/2$}{\rule{1.6cm}{0.4pt}}};
  \node[beablau,below] at (1.5,0) {\lsgoder{$v\cdot\Delta t/2$}{\rule[-3mm]{1.6cm}{0.4pt}}};
  \node[zoeorange,right] at (3,1.1) {\lsgoder{$c\cdot\Delta t_0/2$}{\rule{1.6cm}{0.4pt}}};
  \node[font=\scriptsize,below] at (6,0) {Tick};
\end{tikzpicture}\\
{\scriptsize Ein Tick von Zoes Uhr in Beas System}
\end{minipage}\hfill
\begin{minipage}[t]{0.57\linewidth}
\vspace{0pt}\lueckentext
Derselbe Tick dauert für \bea{Bea} $\Delta t$, für \zoe{Zoe} $\Delta t_0$.\\
\nr{1} \textbf{Seiten} beschriften (halbe Tick-Dauer, im Bild).\\
\nr{2} \textbf{Pythagoras:}\\ \lsgrest{$\left(c\,\frac{\Delta t}{2}\right)^2 = \left(v\,\frac{\Delta t}{2}\right)^2 + \left(c\,\frac{\Delta t_0}{2}\right)^2$}
\nr{3} \textbf{Auflösen} nach $\Delta t$:\\ \lsgrest{$c^2\Delta t^2 - v^2\Delta t^2 = c^2\Delta t_0^2$}
\lsgrest{$\Delta t^2\,(1 - v^2/c^2) = \Delta t_0^2$}
$\Delta t = $ \lsg[5cm]{$\Delta t_0 / \sqrt{1 - v^2/c^2}$}
\end{minipage}\\[1mm]
\lueckentext
Prüfen mit $v = 0{,}6\,c$: \lsgrest{$\sqrt{1 - 0{,}36} = 0{,}8$; \quad $\Delta t = 10$ ns : 0,8 = 12,5 ns \checkmark}
\end{aufgabe}

\newpage
\begin{aufgabe}[Gilt das nur für Lichtuhren? \normalfont(Argumentiere mit Postulat 1)]{4}
\lsglinien{3}{Nein. Ginge Zoes Armbanduhr anders als ihre Lichtuhr, könnte sie im geschlossenen Zug feststellen, dass sie fährt. Das verbietet Postulat 1. Also laufen alle Uhren (auch Puls, Zerfälle) im Gleichschritt.}
\end{aufgabe}

\begin{merke}[Zeitdilatation]
\lueckentext
Die \textbf{Eigenzeit} $\Delta t_0$ misst, für wen beide Ereignisse \lsg[5.25cm]{am selben Ort} stattfinden.\\
In jedem anderen Inertialsystem liegt zwischen den Ereignissen eine \lsg[3.2cm]{längere} Zeit:\\
$\Delta t = $ \lsg[3cm]{$\gamma\cdot\Delta t_0$} \quad mit dem Lorentzfaktor \quad $\gamma = $ \lsg[4.5cm]{$1/\sqrt{1 - v^2/c^2}$}\\
$\gamma$ ist immer \lsg[1.5cm]{$\geq 1$}. Bewegte Uhren gehen \lsg[3.2cm]{langsamer}. Die Eigenzeit ist die \lsg[3.2cm]{kürzeste} Zeit.
\end{merke}

\begin{aufgabe}[Warum hat das vorher niemand bemerkt? \normalfont($\gamma$-Explorer)]{5}
{\renewcommand{\arraystretch}{1}\small
\begin{tabularx}{\linewidth}{|l|>{\centering\arraybackslash}X|>{\centering\arraybackslash}X|>{\centering\arraybackslash}X|>{\centering\arraybackslash}X|}\hline
 & \textbf{ICE} & \textbf{Flugzeug} & \textbf{ISS} & \textbf{Myon (0,995 c)}\\\hline
$\gamma$ bzw. $\gamma - 1$\schreibzeile & \lsgtext{$1 + 3{,}8\cdot10^{-14}$} & \lsgtext{$1 + 3{,}5\cdot10^{-13}$} & \lsgtext{$1 + 3{,}3\cdot10^{-10}$} & \lsgtext{10,0}\\\hline
Uhr geht nach pro Tag\schreibzeile & \lsgtext{3,3 ns} & \lsgtext{30 ns} & \lsgtext{28 µs} & \lsgtext{(über 21 h)}\\\hline
\end{tabularx}}

\vspace{1mm}\lueckentext
Mein Tipp (ICE): \luecke[2.6cm] \quad $\gamma$ weicht erst ab etwa \lsg[2cm]{0,14 c} um mehr als 1\,\% von 1 ab.
\end{aufgabe}

\begin{aufgabe}[Mit dem Plan rechnen: Raumschiff mit $v = 0{,}8\,c$, Borduhr zeigt 3,0 Jahre]{6}
\lueckentext
\nr{1} \bea{S:} \lsg[4cm]{Erde und Station} \zoe{S$'$:} \lsg[3cm]{Raumschiff} Ereignisse: \lsg[4.6cm]{Abflug und Ankunft}\\
\nr{2} Beide Ereignisse am selben Ort für \lsg[3.5cm]{das Raumschiff}: Eigenzeit misst die \lsg[3cm]{Borduhr}.\\
\nr{3} $\gamma = $ \lsg[5.5cm]{$1/\sqrt{1 - 0{,}64} = 1/0{,}6 = 1{,}67$} \quad $\Delta t = $ \lsg[4.5cm]{$1{,}67\cdot 3{,}0$ a $= 5{,}0$ a}\\
\phantom{\nr{3}} Entfernung der Station für die Erde: \lsgrest{$0{,}8\,c\cdot 5{,}0$ a $= 4{,}0$ Lichtjahre}
\end{aufgabe}

\begin{aufgabe}[Im Experiment geprüft: Atomuhren im Flugzeug (1971)]{7}
{\small Vorhersage / Messung: ostwärts $-40 \pm 23$ ns / $-59 \pm 10$ ns; \quad westwärts $+275 \pm 21$ ns / $+273 \pm 7$ ns.}\\[0.5mm]
\lueckentext
\textbf{a)} Passt das zusammen? \lsgrest{Ja, die Bereiche überlappen jeweils.}
\textbf{b)} Warum geht die Uhr im Ostflug nach, die im Westflug vor?
\lsglinien{2}{Von außen gesehen ist die Uhr im Ostflug am schnellsten (Erddrehung plus Flug), die im Westflug am langsamsten. Die schnellste Uhr geht am langsamsten.}
\end{aufgabe}

\begin{planbox}[Plan-Karte, Stand nach Station 4 \normalfont(ins Logbuch übertragen, Annahmen-Tafel Spalte 4)]
\lueckentext
\textbf{\nr{2} Eigen-System:} beide Ereignisse am \lsg[3cm]{selben Ort} $\to$ Eigenzeit. \quad \textbf{\nr{3} Umrechnen:} Zeit \lsg[1.2cm]{mal} $\gamma$.
\end{planbox}
\end{document}
)
%  Lösungseinträge belegen nie mehr Platz als die Lücke im Blatt.
%
%  Antwortplatz:
%   \luecke[l] / \lsg[l]{Text}      Lücke (Standard 6,75 cm, Fachbegriff 5,25 cm, Zahl 2,25 cm)
%   \restzeile / \lsgrest{Text}     Lücke bis Zeilenende      \lueckentext  Absatz mit Lücken
%   \linien{n} / \lsglinien{n}{T}   Schreiblinien 8 mm        \schreibzeile Tabellenzelle 8 mm
%   \lsgtext{…}                     nur in der Lösung         \lsgoder{Lösung}{Blatt}
%   \kariert{n} / \lsgkariert{n}{TikZ}, \karobox{b}{h} / \lsgkarobox{b}{h}{TikZ}
%   \kreuz{1}                       Ankreuzfeld (1 = in der Lösung angekreuzt)
%   \nr{1}                          Schrittnummer ① wie auf der Website
%  Feste Farben wie auf der Website:
%   \bea{…} blau (Bahnsteig-System S), \zoe{…} orange (Zug-System S′), \licht{…} gelb
%   \weblink{Abschnitt}             grauer Verweis auf den Abschnitt der Stationsseite
% =====================================================================
\usepackage[a4paper,left=16mm,right=16mm,top=23mm,bottom=13mm,headheight=12mm,headsep=3mm,footskip=6mm]{geometry}
\usepackage{iftex}
\ifPDFTeX
  \usepackage[T1]{fontenc}
  \usepackage[utf8]{inputenc}
  \usepackage{helvet}
  \renewcommand{\familydefault}{\sfdefault}
\else
  \usepackage{fontspec}
  \setmainfont{texgyreheros}[Extension=.otf,UprightFont=*-regular,BoldFont=*-bold,ItalicFont=*-italic,BoldItalicFont=*-bolditalic]
\fi
\usepackage[ngerman,shorthands=off]{babel}
\usepackage{amsmath,amssymb}
\usepackage{xcolor,graphicx,tabularx,array,enumitem,fancyhdr,tikz,multicol,calc,etoolbox,pifont,xspace,colortbl}
\usepackage[most]{tcolorbox}
\usetikzlibrary{arrows.meta,positioning,calc}
\usepackage[hidelinks]{hyperref}
\setlength{\parindent}{0pt}
\setlength{\parskip}{3pt}
\setlist{nosep,leftmargin=*}

% ---------- Lösungsschalter ----------
\newif\ifloesung
\ifdefined\mitloesung\loesungtrue\fi
\newcommand{\lsgfarbe}{\color{red!75!black}}

% ---------- Farben wie auf der Website ----------
\definecolor{beablau}{HTML}{2563EB}\definecolor{zoeorange}{HTML}{EA580C}\definecolor{lichtgelb}{HTML}{CA8A04}
\definecolor{relgruen}{HTML}{0F766E}
\newcommand{\bea}[1]{{\color{beablau}\bfseries #1}}
\newcommand{\zoe}[1]{{\color{zoeorange}\bfseries #1}}
\newcommand{\licht}[1]{{\color{lichtgelb}\bfseries #1}}

% ---------- Kopf- und Fußzeile ----------
\newcommand{\lehrkraft}{Hau}
\newcommand{\themenbereich}{3.2 Spezielle Relativitätstheorie}
\newcommand{\abnummer}{}
\pagestyle{fancy}\fancyhf{}
\renewcommand{\headrulewidth}{0.4pt}
\fancyhead[L]{\small Klasse 11\\Datum:}
\fancyhead[C]{\small Themenbereich:\\\themenbereich}
\fancyhead[R]{\small \lehrkraft\ifloesung\\{\lsgfarbe\bfseries Lösung}\fi}
\fancyfoot[C]{\scriptsize edu-mrh.de\,·\,Physik 11\,·\,Blatt \abnummer\ifloesung\ (Lösung)\fi\,·\,Seite \thepage}
\newcommand{\abtitel}[2]{\renewcommand{\abnummer}{#1}\begin{center}\Large\bfseries #1\quad\underline{#2}\end{center}\vspace{-2mm}}
\newcommand{\leitfrage}[1]{\begin{center}\vspace{-1.5mm}\itshape #1\end{center}\vspace{-1.5mm}}

% ---------- Boxen ----------
\newenvironment{aufgabe}[2][]{\begin{tcolorbox}[enhanced,breakable,colback=white,colframe=black,boxrule=0.7pt,arc=3mm,left=2mm,right=2mm,top=1mm,bottom=1.2mm,before skip=2mm,after skip=2mm]\textbf{Aufgabe #2}\ifstrempty{#1}{}{ \textit{#1}}\par\vspace{0.3mm}}{\end{tcolorbox}}
\newtcolorbox{merke}[1][Merke]{enhanced,breakable,colback=green!5,colframe=green!40!black,boxrule=0.8pt,arc=2mm,left=2mm,right=2mm,top=1mm,bottom=1.2mm,before skip=2mm,after skip=2mm,title={#1},fonttitle=\bfseries,coltitle=white,toptitle=0.3mm,bottomtitle=0.3mm}
\newtcolorbox{hinweis}{enhanced,breakable,colback=yellow!12,colframe=orange!70!black,boxrule=0.6pt,arc=2mm,left=2mm,right=2mm,top=1mm,bottom=1mm,before skip=2mm,after skip=2mm}
\newtcolorbox{planbox}[1][Plan-Karte]{enhanced,breakable,colback=relgruen!5,colframe=relgruen,boxrule=0.9pt,arc=2mm,left=2mm,right=2mm,top=1mm,bottom=1.2mm,before skip=2mm,after skip=2mm,title={#1},fonttitle=\bfseries,coltitle=white,toptitle=0.3mm,bottomtitle=0.3mm}
\newcommand{\weblink}[1]{{\small\color{gray}\textit{Website: #1}}}

% ---------- Antwortplatz (wie informatik/10/java, Lücken x 1,5) ----------
\newlength{\zeilenabstand}\setlength{\zeilenabstand}{8mm}
\newlength{\karo}\setlength{\karo}{5mm}
\tikzset{karo/.style={draw=gray!50,line width=0.3pt},schreiblinie/.style={draw=black!65,line width=0.35pt}}
\newlength{\lsgbreite}
\newcommand{\lsg}[2][6.75cm]{\leavevmode\ifloesung\settowidth{\lsgbreite}{\,#2\,}%
  \ifdim\lsgbreite>#1\relax\rlap{\raisebox{0.8pt}{\resizebox{#1}{!}{\lsgfarbe\,#2\,}}}%
  \else\rlap{\raisebox{0.8pt}{\lsgfarbe\,#2}}\fi\fi\rule[-0.6ex]{#1}{0.4pt}\xspace}
\newcommand{\luecke}[1][6.75cm]{\lsg[#1]{}}
\newcommand{\lsgrest}[1]{\leavevmode\unskip\hspace{1.5mm}\ifloesung\rlap{\raisebox{0.8pt}{\lsgfarbe\,#1}}\fi\leaders\hbox{\rule[-0.6ex]{0.5pt}{0.4pt}}\hfill\null\par}
\newcommand{\restzeile}{\lsgrest{}}
\newcommand{\lsgtext}[1]{\ifloesung\leavevmode{\lsgfarbe#1}\fi}
\newcommand{\lsgoder}[2]{\ifloesung\leavevmode#1\else#2\fi}
\newcommand{\lueckentext}{\linespread{1.5}\selectfont\raggedright}
\newcommand{\schreibzeile}{\rule[-2.8mm]{0pt}{8mm}}
\newcommand{\nr}[1]{\ding{\numexpr171+#1\relax}}
\newcommand{\lsglinien}[2]{\par\vspace{0.5mm}\noindent\begin{tikzpicture}
  \path (0,0) (\linewidth-0.4pt,-#1\zeilenabstand-1.5mm);
  \foreach \i in {1,...,#1}{\draw[schreiblinie] (0.2pt,-\i\zeilenabstand) -- (\linewidth-0.2pt,-\i\zeilenabstand);}
  \ifloesung\node[overlay,anchor=north west,inner sep=0pt] at (1mm,-\zeilenabstand+1pt+5mm) {\parbox[t]{\linewidth-2mm}{\lsgfarbe\raggedright\setlength{\baselineskip}{\zeilenabstand}\rule{0pt}{5mm}#2}};\fi
  \end{tikzpicture}\par}
\newcommand{\linien}[1]{\lsglinien{#1}{}}
\newcommand{\karoraster}[3]{\begin{tikzpicture}\draw[karo] (0,0) grid[step=\karo] (#1\karo,#2\karo);
  \ifloesung\begin{scope}[overlay,shift={(0,#2\karo)},red!75!black,text=red!75!black,line width=0.8pt]#3\end{scope}\fi\end{tikzpicture}}
\newcommand{\lsgkariert}[2]{\par\vspace{1.5mm}\noindent\pgfmathtruncatemacro{\karoX}{floor(\linewidth/\karo)}%
  \makebox[\linewidth]{\karoraster{\karoX}{#1}{#2}}\par\vspace{1mm}}
\newcommand{\kariert}[1]{\lsgkariert{#1}{}}
\newcommand{\lsgkarobox}[3]{\pgfmathtruncatemacro{\karoX}{round((#1)/\karo)}\pgfmathtruncatemacro{\karoY}{round((#2)/\karo)}\karoraster{\karoX}{\karoY}{#3}}
\newcommand{\karobox}[2]{\lsgkarobox{#1}{#2}{}}
\newcommand{\kreuz}[1]{\leavevmode\ifloesung\ifnum#1=1 {\lsgfarbe$\boxtimes$}\else$\square$\fi\else$\square$\fi}

% ---------- Zeichnungen: Strichfigur (Füße im Ursprung des scope) ----------
%  \figur{Farbe}{Name} in einer tikzpicture, vorher mit \begin{scope}[shift={(x,y)}] platzieren
\newcommand{\figur}[2]{\draw[#1,line width=1pt,line cap=round] (0,0.62) -- (0,0.28) (0,0.28) -- (-0.12,0) (0,0.28) -- (0.12,0) (-0.15,0.48) -- (0.15,0.48);
  \fill[#1] (0,0.74) circle (0.11); \node[#1,font=\scriptsize\bfseries,above] at (0,0.82) {#2};}

% Blatt:   xelatex ab-2-licht.tex
% Lösung:  xelatex ab-2-licht-lsg.tex
% Progression: Vorhersage Taschenlampe -> Äther als Analogieschluss -> Wettschwimmen (Vorhersage, Rechnung, Dreieck)
%              -> Michelson-Morley (Erwartung/Messung) -> Hypothesen-Gericht -> Postulate -> Schlüsselgedanke
\begin{document}
\abtitel{2}{Das Problem mit dem Licht}
\leitfrage{Licht hat die Geschwindigkeit c = 300\,000 km/s. Aber relativ wozu?}

\begin{aufgabe}[Vorhersage: Die Taschenlampe im Zug]{1}
\zoe{Zoes} Zug fährt mit $100\,000$ km/s. Sie leuchtet nach vorn und misst für das Licht 300\,000 km/s. \bea{Bea} misst:\\[0.5mm]
\kreuz{0} 400\,000 km/s ($c + v$) \quad \kreuz{1} 300\,000 km/s \quad \kreuz{0} 200\,000 km/s ($c - v$) \hfill {\small (Auflösung in Aufgabe 5)}\\[1mm]
\lueckentext
Um 1880: Schall braucht Luft, also braucht Licht auch ein Medium, den \lsg[3.5cm]{Äther}. Die Erde bewegt sich mit 30 km/s.\\
Dann müsste man auf der Erde bemerken:\\ \lsgrest{einen „Ätherwind“: Licht wäre je nach Richtung verschieden schnell.}
\end{aufgabe}

\begin{aufgabe}[Wettschwimmen: quer gegen längs]{2}
\begin{minipage}[t]{0.56\linewidth}
\vspace{0pt}
Beide schwimmen mit $c = 2$ m/s relativ zum Wasser, Strömung $v = 1{,}2$ m/s, jede Strecke 100 m hin und 100 m zurück.\\[1mm]
Mein Tipp: Zuerst zurück ist \quad \kreuz{1} \zoe{A (quer)} \quad \kreuz{0} \bea{B (längs)} \quad \kreuz{0} beide\\[1mm]
\lueckentext
\textbf{a)} Simulation ($v = 0{,}6\,c$): A \lsg[1.5cm]{5,00 s}, B \lsg[1.5cm]{6,25 s}, ohne Strömung \lsg[1.5cm]{4,00 s}\\
\textbf{b)} \bea{B} flussabwärts: Tempo \lsg[2cm]{3,2 m/s}, Zeit \lsg[2cm]{31,25 s}\\
\phantom{\textbf{b)}} flussaufwärts: Tempo \lsg[2cm]{0,8 m/s}, Zeit \lsg[2cm]{125 s}\\
\phantom{\textbf{b)}} zusammen: $t_B = $ \lsg[3cm]{156,25 s}
\end{minipage}\hfill
\begin{minipage}[t]{0.4\linewidth}
\vspace{0pt}\centering
\begin{tikzpicture}[x=1cm,y=1cm,font=\small]
  \fill[cyan!10] (0,0) rectangle (6.2,3.2);
  \draw[-{Stealth},cyan!60!black,line width=1pt] (4.2,2.8) -- (5.8,2.8) node[midway,below] {Strömung $v$};
  \fill (0.8,0.3) circle (2pt) node[below] {Start}; \fill (0.8,2.9) circle (2pt); \fill (3.4,0.3) circle (2pt);
  \draw[zoeorange,line width=1pt,{Stealth}-{Stealth}] (0.8,0.45) -- (0.8,2.75) node[midway,left] {A};
  \draw[beablau,line width=1pt,{Stealth}-{Stealth}] (0.95,0.3) -- (3.25,0.3) node[midway,above] {B};
  \draw[beablau,line width=1.2pt] (3.9,0.9) -- (5.1,0.9) node[midway,below] {$v$};
  \draw[zoeorange,line width=1.2pt] (5.1,0.9) -- (5.1,2.5);
  \draw[lichtgelb,line width=1.2pt] (3.9,0.9) -- (5.1,2.5) node[midway,left] {$c$};
  \node[zoeorange,right] at (5.1,1.7) {?};
\end{tikzpicture}
\end{minipage}\\[1mm]
\lueckentext
\textbf{c)} \zoe{A} muss schräg gegen die Strömung halten. Quer-Tempo (Pythagoras im Dreieck):\\ \lsgrest{$\sqrt{c^2 - v^2} = \sqrt{4 - 1{,}44}$ m/s $= 1{,}6$ m/s}
\phantom{\textbf{c)}} $t_A = $ \lsg[4.5cm]{200 m : 1,6 m/s = 125 s} \quad ohne Strömung beide: \lsg[2.25cm]{100 s}\\
\textbf{d)} Sind die beiden Zeiten verschieden, dann weiß ich: \lsgrest{Das Wasser strömt (und kann $v$ ausrechnen).}
\end{aufgabe}

\begin{aufgabe}[Das Experiment von Michelson und Morley (1887): Wettrennen mit Licht]{3}
{\renewcommand{\arraystretch}{1}\small
\begin{tabularx}{\linewidth}{|l|>{\raggedright\arraybackslash}X|>{\raggedright\arraybackslash}X|}\hline
 & \textbf{Vorhersage mit Ätherwind} & \textbf{Messung}\\\hline
Streifenmuster beim Drehen\rule[-6.5mm]{0pt}{14mm} & \lsgtext{verschiebt sich um etwa 0,4 Streifenbreiten (je 90° wechselt der Vorsprung)} & \lsgtext{bleibt stehen: Verschiebung unter 0,01, zu jeder Tages- und Jahreszeit}\\\hline
\end{tabularx}}

\vspace{1mm}\lueckentext
\textbf{c)} Übersetzt ins Wettschwimmen heißt das:\\ \lsgrest{A und B kommen immer gleichzeitig an: Der „Fluss“ strömt nicht.}
\end{aufgabe}

\newpage
\begin{aufgabe}[Hypothesen-Gericht: Welches Beweisstück widerlegt den Rettungsversuch?]{4}
\textbf{J} Jahreszeiten \quad \textbf{S} schräge Teleskope \quad \textbf{D} Doppelsterne und schnelle Pionen \quad \textbf{G} Genauigkeit\\[1mm]
{\renewcommand{\arraystretch}{1}\small
\begin{tabularx}{\linewidth}{|>{\raggedright\arraybackslash}p{7.2cm}|c|>{\raggedright\arraybackslash}X|}\hline
\textbf{Rettungsversuch} & \textbf{Nr.} & \textbf{Warum er scheitert}\\\hline
„Das Gerät war zu ungenau.“\rule[-4.5mm]{0pt}{11mm} & \lsgtext{G} & \lsgtext{Es hätte schon 1/40 des Effekts gezeigt.}\\\hline
„Die Erde stand zufällig still im Äther.“\rule[-4.5mm]{0pt}{11mm} & \lsgtext{J} & \lsgtext{Ein halbes Jahr später fliegt sie entgegengesetzt, das Ergebnis bleibt null.}\\\hline
„Die Erde reißt den Äther mit.“\rule[-4.5mm]{0pt}{11mm} & \lsgtext{S} & \lsgtext{Dann müsste man Teleskope nicht neigen.}\\\hline
„Licht verhält sich wie ein Ball ($c + v$).“\rule[-4.5mm]{0pt}{11mm} & \lsgtext{D} & \lsgtext{Licht bewegter Quellen kommt trotzdem mit c an.}\\\hline
\end{tabularx}}
\end{aufgabe}

\begin{merke}[Einsteins zwei Postulate (1905)]
\textbf{Postulat 1 (Relativitätsprinzip):}
\lsglinien{2}{In allen Inertialsystemen gelten dieselben Naturgesetze (nicht nur die der Mechanik). Kein Experiment zeichnet ein Inertialsystem als „ruhend“ aus.}
{\small Beobachtung dahinter:} \lsgrest{Galileis Schiff; kein Ätherwind messbar}
\textbf{Postulat 2 (Konstanz der Lichtgeschwindigkeit):}
\lsglinien{2}{Licht breitet sich im Vakuum in jedem Inertialsystem mit derselben Geschwindigkeit c aus, egal wie sich Quelle oder Beobachter bewegen.}
{\small Beobachtung dahinter:} \lsgrest{Michelson-Morley, Doppelsterne, schnelle Pionen}
\lueckentext
$c = $ \lsg[4.5cm]{299\,792\,458 m/s} $\approx$ \lsg[3.5cm]{$3{,}00\cdot10^8$ m/s}. Ein Äther wird \lsg[3.5cm]{nicht gebraucht}.
\end{merke}

\begin{aufgabe}[Der Konflikt: Irgendetwas muss nachgeben]{5}
Auflösung zu Aufgabe 1: Bea misst \lsg[3.5cm]{300\,000 km/s}. Die Regel $u = u' + v$ versagt für Licht.\\[1mm]
\centerline{\large Geschwindigkeit $=$ Strecke $:$ Zeit}\par\vspace{1mm}
Bea und Zoe messen für dasselbe Licht dieselbe Geschwindigkeit, obwohl sie sich gegeneinander bewegen. Mein Schluss:
\lsglinien{3}{Dann können Strecken und Zeiten nicht für alle gleich sein. Bea und Zoe müssen für denselben Vorgang verschiedene Zeiten und Längen messen, gerade so, dass bei beiden c herauskommt.}
\end{aufgabe}

\begin{planbox}[Plan-Karte, Stand nach Station 2 \normalfont(ins Logbuch übertragen, Annahmen-Tafel Spalte 2 ausfüllen)]
\lueckentext
\textbf{$\bigstar$} P1 und P2 (Merkkasten). \quad \textbf{Bild:} das Dreieck aus Aufgabe 2 mit $c$, $v$ und \lsg[3cm]{$\sqrt{c^2 - v^2}$}.\\
Nächste Fragen: Station 3 \lsg[2.5cm]{wann?}, Station 4 \lsg[2.5cm]{wie lange?}, Station 5 \lsg[2.5cm]{wie lang?}
\end{planbox}
\end{document}
